\section{Near-Collision, Photon Production, and Entanglement}
\begin{quote}\small
\noindent\textbf{Status.} \emph{Legacy physical-regime block inside the active main body.}\par
\noindent\textbf{Block function.} This section records how the same closure architecture is read in the near-collision and entanglement regime, together with the local role of the photon readout.
\end{quote}
\label{sec:entanglement}
%% ============================================================

\begin{proposition}[Compensator as near-collision in $K_3$]\label{prop:compensator}
The Heisenberg compensator $C(f)=\tfrac{1}{2}(f+Jf)$ is the projection
onto $K_3=e_1 e_2$: it takes matter $f$ (in $K_1$) and its tachyon image
$Jf$ (in $K_2$), and produces the photon $C(f)\in K_3$.
The quaternion identity $e_3=e_1 e_2$ ensures this is always orthogonal
to both inputs.
\end{proposition}

\begin{theorem}[Entanglement as shared proper time]
\label{thm:entanglement-shared-time}
Two photons $\phi_1,\phi_2\in K_3$ are entangled iff they share the same
pre-image: $\phi_1=\phi_2=C(f)$ for some $f\in K_1$.
The entanglement correlation is the invariant proper time
$\tau_n(f)=\tau_n(Jf)$ fixed at the moment of the near-collision.
No signal is transmitted at measurement: the invariant $\tau_n$ was
established at the near-collision and is merely revealed.
\end{theorem}

\begin{proof}
$\tau_n(Jf)^2 = t_n^2 - |-\zminus|^2 = t_n^2 - |\zminus|^2 = \tau_n(f)^2$
by Theorem~\ref{thm:lorentz}.
The shared value $\tau_n$ is a Lorentz invariant at every stratum,
hence frame-independent.
\end{proof}

%% ============================================================

